\documentclass[
	% -- opções da classe memoir --
	12pt,				    % tamanho da fonte
	openright,			    % capítulos começam em pág ímpar
	oneside,			    % para impressão só no anverso. Oposto a twoside
	a4paper,			    % tamanho do papel.
    % -- opções do pacote abntex2 --
    chapter=TITLE,         % Títulos em maiúsculas
    sumario=tradicional,    % Sumário padrão memoir
    % -- opções do pacote babel --
	english,			    % idioma adicional para hifenização
	brazil,				    % o último idioma é o principal do documento
	]{abntex2}              % Personaliza a capa. Precisa do arquivo ufv.cls para funcionar.



% Pacotes fundamentais
\usepackage{abntex2-UFV}        % Personalização para a Universidade Federal de Viçosa
\usepackage{newtxtext,newtxmath}
\usepackage[T1]{fontenc}		% Selecao de codigos de fonte de saída
\usepackage[utf8]{inputenc}		% Codificacao do documento (conversão automática dos acentos)
\usepackage{indentfirst}		% Indenta o primeiro parágrafo de cada seção.
\usepackage{graphicx}			% Inclusão de gráficos
\usepackage{booktabs}           % \toprule, \midrule e \bottomrule para tabelas
% \usepackage[none]{hyphenat}
% Sistema autor-data com títulos nas referências em negrito
\usepackage[abnt-emphasize=bf]{abntex2cite}	
\usepackage{tabularx}
\usepackage{multirow}
\usepackage{float}
\usepackage[portuguese,onelanguage,ruled,vlined]{algorithm2e}

% ---
% CONFIGURAÇÕES DE PACOTES
% ---

% Informações de dados para CAPA e FOLHA DE ROSTO
\titulo{Título do relatório}
\autor{Nome do Autor}
\local{Florestal}
\data{Setembro 2021}
\orientador{Nome do Professor}
\instituicao{Universidade Federal de Viçosa}
\instituto{INSTITUTO DE CIÊNCIAS EXATAS E TECNOLÓGICAS}
\campus{\emph{Campus} de Florestal}
\curso{Bacharelado em Ciência da Computação}

% O preambulo deve conter o tipo do trabalho, o objetivo,
% o nome da instituição e a área de concentração
\preambulo{Relatório Final, referente ao período de MM/AAAA a MM/AAAA, apresentado à Universidade Federal de Viçosa, como parte das exigências do PIBIC/CNPq.}
% ---

% ---
% Configurações de aparência do PDF final

% informações para o arquivo pdf de saída
% Interessante alterar a cor dos links para preto(black)
% para imprimir
\makeatletter
\hypersetup{
        % metadados
		pdftitle={\@title},
		pdfauthor={\@author},
    	pdfsubject={\imprimirpreambulo},
	    pdfcreator={LaTeX with abnTeX2},
		colorlinks=true,   % false: links em frame; true: links coloridos
    	linkcolor=black,    % cor dos links no documento
    	citecolor=blue,    % cor dos links para a bibliografia
    	filecolor=magenta, % cor dos links para arquivos
		urlcolor=blue,     % cor dos links para sites
		bookmarksdepth=4   % profundidade do sumário do PDF
}
\makeatother
% ---

\begin{document}
% Retira espaço extra obsoleto entre as frases.
\frenchspacing

% ----------------------------------------------------------
% ELEMENTOS PRÉ-TEXTUAIS
% ----------------------------------------------------------
\pretextual

% Folha de rosto
\imprimirfolhaderosto

% ----------------------------------------------------------
% Resumo
% ----------------------------------------------------------
\input{sections/0-resumo}


% % inserir lista de ilustrações
% \pdfbookmark[0]{\listfigurename}{lof}
% \listoffigures*
% \cleardoublepage
% % ---

% % inserir lista de tabelas
% \pdfbookmark[0]{\listtablename}{lot}
% \listoftables*
% \cleardoublepage
% % ---

% Lista de siglas e abreviaturas (opcional)
% sintaxe: \item [sigla] Descrição da sigla

%\begin{siglas}
%\item[ABNT] Absurdas Normas Técnicas
%\item[UFV] Universidade Federal de Viçosa
%\item[CAF] \emph{Campus} de Florestal
%\end{siglas}

% Lista de símbolos (opcional)
% sintaxe: \item [simbolo] Descrição do símbolo

%\begin{simbolos}
%\item[$\infty$ ] Infinito
%\end{simbolos}


% inserir o sumario
\pdfbookmark[0]{\contentsname}{toc}
\tableofcontents*
\cleardoublepage
% ---



% ----------------------------------------------------------
% ELEMENTOS TEXTUAIS
% ----------------------------------------------------------
\textual

\gdef\clearforchapter{}
% ----------------------------------------------------------
% Introdução
% ----------------------------------------------------------
\input{sections/1-introducao} \\

% ----------------------------------------------------------
% Revisão da Literatura
% ----------------------------------------------------------
\input{sections/2-relacionados} \\

% ----------------------------------------------------------
% Objetivos
% ----------------------------------------------------------
\input{sections/3-objetivos} \\

% ----------------------------------------------------------
% Material e Métodos
% ----------------------------------------------------------
\input{sections/4-material-e-metodos} \\

% ----------------------------------------------------------
% Resultados
% ----------------------------------------------------------
\input{sections/5-resultados} \\

% ----------------------------------------------------------
% Conclusões
% ----------------------------------------------------------
\input{sections/6-conclusoes} \\


% ----------------------------------------------------------
% ELEMENTOS PÓS-TEXTUAIS
% ----------------------------------------------------------
\cleardoublepage
\postextual

% Referências bibliográficas

\bibliography{referencias}

% Caso sejam necessários apêndices ou anexos em seu documento
% Use os ambientes abaixo

%% Apêndices
%
%% Inicia os apêndices
%\begin{apendicesenv}
%
%% Imprime uma página indicando o início dos apêndices
%\partapendices
%
%\chapter{Primeiro Apêndice}
%
%\chapter{Segundo Apêndice}
%
%\end{apendicesenv}
%
%
%% ----------------------------------------------------------
%% Anexos
%% ----------------------------------------------------------
%\begin{anexosenv}
%
%% Imprime uma página indicando o início dos anexos
%\partanexos
%
%\chapter{Primeiro Anexo}
%\lipsum[30]
%
%\chapter{Segundo Anexo}
%\lipsum[31]
%
%\end{anexosenv}

\end{document}
